
\vspace{-0.5em}
\section{如何在门诊（OPD）与半门诊（Semi-OPD）之间进行选择？}
\vspace{-0.5em}
\label{sec:when_why}

如第~\ref{sec:experiments} 节和图~\ref{fig:allpairs} 所示，OPD 与 Semi-OPD 的相对性能与教师-学生重叠比例高度相关。因此，该指标可作为判断何时使用 OPD 或 Semi-OPD 的指示性依据。下文将更正式地描述这一相关系数。

\noindent \textbf{当教师-学生对的重叠比例较高时，需要使用OPD；否则，Semi-OPD更优且更快。} 我们在图~\ref{fig:overlap} 中计算了Semi-OPD相对于OPD的增益（以两者学习曲线之间的面积除以共享训练时段进行规范化的结果）与每对教师-学生重叠比例之间的相关系数。这两个量具有很强的相关性（17对中的皮尔逊相关系数为$r=-0.73$）：初始重叠比例较低的对从Semi-OPD中获益更大，而只有重叠比例较高的对才需要使用OPD。

该模式表明，当学生模型与教师模型存在显著差异时，反复从不断演化的学生策略中采样会导致在线回放质量下降（若干示例见附录~\ref{app:qualitative}），并削弱监督信号；而重用初始学生模型的回放数据可提供更稳定的蒸馏信号。我们将在第~\ref{sec:why} 节深入探讨其潜在原因。



\textbf{关键见解。}
上述现象为OPD的实际应用与研究提供了有益启示：(1) 针对特定的目标教师-学生对和目标领域提示，应根据其重叠比例来决定是否使用OPD或Semi-OPD，而非盲目采用OPD；(2) 未来在语言模型后训练流水线中研究改进的OPD方案时，需谨慎选择作为实验情景的教师-学生对；否则，Semi-OPD可作为强有力的基准。



\textbf{在教师生成的回放序列上进行SFT能否缩小差距？}
我们还进行了两组实验，其中学生模型在教师生成的回放序列上进行一轮SFT，随后进行OPD或半监督OPD。这与Lightning OPD~\citep{wu2026lightning}的设置一致。如图~\ref{fig:sft}所示，我们观察到：(1) 在教师回放序列上进行SFT略微提高了教师-学生重叠率，但半监督OPD在两组实验中仍优于OPD，尽管优势幅度较小；(2) 我们的半监督OPD在不包含额外的教师回放序列SFT阶段的情况下，性能反而优于包含该阶段的半监督OPD，表明在教师回放序列上进行SFT是不必要的。
此外我们注意到，尽管~\citep{wu2026lightning}将固定回放序列主要视为OPD的一种高效近似，但我们的结果表明，当初始教师-学生重叠率较低时，冻结回放缓存反而能提升蒸馏效果。


